% Mathematical foundations of the dormant-timebomb project.
% Companion to docs/mathematical_foundations.md; the two are kept in step by hand.
%
%   pdflatex mathematical_foundations.tex
%
\documentclass[11pt,a4paper]{article}

\usepackage[utf8]{inputenc}
\usepackage[T1]{fontenc}
\usepackage{amsmath,amssymb}
\usepackage{booktabs}
\usepackage{geometry}
\usepackage[hidelinks]{hyperref}
\usepackage{microtype}

\geometry{margin=1in}

% \code{...} survives underscores and other specials in identifiers and paths
\newcommand{\code}[1]{\texttt{\detokenize{#1}}}
\newcommand{\relu}{\operatorname{relu}}
\newcommand{\softmax}{\operatorname{softmax}}
\newcommand{\logit}{\operatorname{logit}}
\newcommand{\sigmoid}{\operatorname{sigmoid}}
\newcommand{\Bin}{\operatorname{Binomial}}
\newcommand{\CE}{\mathrm{CE}}
\newcommand{\PPL}{\mathrm{PPL}}
\DeclareMathOperator*{\argmin}{arg\,min}

\title{Mathematical foundations}
\author{Dormant timebombs: adversarial attacks against collapsed code models}
\date{\today}

\begin{document}
\maketitle

Every formula this project rests on, in one place, with a pointer to the code that implements it.
Written for the authors and anyone joining the project: it assumes transformers, cross-entropy and
gradient-based discrete search, and assumes nothing about this repository.

Three threads run through it. \textbf{Collapse} is a dynamical system on a trajectory of models.
\textbf{Surrogates} approximate a point on that trajectory the attacker cannot observe. The
\textbf{attack} is a constrained discrete optimization whose feasibility is governed by a measurable
scalar, the security margin. The last two sections cover the poisoning line of work and a symbol
table.

\tableofcontents

\section{Setup and notation}

A trajectory of models $M_0, M_1, \dots, M_G$ is produced by recursive training. $M_{\mathrm{base}}$
is the pristine pretrained model the trajectory starts from.

\begin{center}
\begin{tabular}{ll}
\toprule
symbol & meaning \\
\midrule
$M_{\mathrm{base}}$ & the pristine base model, never fine-tuned here \\
$M_g$ & the checkpoint after $g$ rounds of recursive training \\
$D_g$ & the corpus generated by $M_g$ \\
$\theta_g$ & the parameters of $M_g$ \\
$l_M(x)$ & the logit vector $M$ assigns at context $x$ \\
$p_M(\cdot \mid x)$ & the next-token distribution, $\softmax(l_M(x))$ \\
$s$ & the adversarial suffix, a token sequence \\
$w$, $r$ & the wrong (insecure) and right (secure) reference implementations \\
$V$, $T$ & vocabulary size, sampling temperature \\
\bottomrule
\end{tabular}
\end{center}

The \textbf{per-token cross-entropy} of a target string $y$ under model $M$, given a prompt built
with suffix $s$, is the quantity every loss below is written in:
\begin{equation}
  \CE_M(y \mid s) \;=\; \frac{1}{|y|} \sum_{t} -\log p_M\!\left( y_t \;\middle|\; \mathrm{prompt}(s),\, y_{<t} \right).
  \label{eq:ce}
\end{equation}
It is a \emph{mean} over the target's tokens, so $|y| \cdot \CE_M(y \mid s) = -\log p_M(y \mid
\mathrm{prompt}(s))$ is the sequence log-likelihood under teacher forcing. Implemented in
\code{TargetModel.target_losses} (\href{../run_attack.py}{\code{run_attack.py}}).

\section{The collapse process}

\subsection{The two recursions}

Generation $g$ trains on what generation $g-1$ wrote. This is the \textbf{data recursion} and it is
the collapse mechanism:
\begin{align}
  D_g &\;\sim\; M_g(\,\cdot \mid \text{instructions}\,), \qquad \text{sampled at } (T, \text{top-}p, \text{top-}k) \\
  M_{g+1} &\;=\; \mathrm{finetune}\!\left( \mathrm{init}_g,\; D_g \right).
\end{align}

The \textbf{weight recursion} is a separate choice. By default $\mathrm{init}_g =
M_{\mathrm{base}}$: each generation trains a new adapter on the pristine base, the ``replace''
setting of the model-collapse literature, so only the data recurs. With \code{--recursive_init},
the ablation, $\mathrm{init}_g = M_g$, so one LoRA adapter keeps being optimized and the number of
optimizer steps applied to it grows linearly in $g$; at a high real-data fraction those steps are
mostly the same human rows again and again, a memorisation route to degradation that
self-training does not cause. The two produce different curves from the same data, which is why
both are in the artifact names. Every result measured before 2026-10-15 is on recursive lineages.

\subsection{The data mixture}

\code{--real_data_fraction} interpolates the training corpus between self-generated and human text
while holding the corpus \emph{size} fixed, so the number of optimizer steps per generation is
constant and the collapse curve stays attributable to composition rather than to budget:
\begin{equation}
  \mathrm{corpus}_g \;=\; (1-f)\cdot \mathrm{sample}(D_{g-1}) \;+\; f \cdot \mathrm{sample}(H),
  \qquad f \in [0,1],\; g \ge 1 .
  \label{eq:mixture}
\end{equation}

$f = 0$ is the pure replace regime. $f = 1$ is the \textbf{non-collapsing control}: same generation
count, same seeds, same step count, no self-generated data at all, which is what separates
``self-training moved the model'' from ``iterated fine-tuning moved the model''. Generation 0 is
never mixed, since it trains on human data by definition. Implemented in \code{mix_real_data}
(\href{../run_baseline.py}{\code{run_baseline.py}}).

\subsection{Why truncated sampling drives collapse}
\label{sec:sharpening}

Sampling at temperature $T$ with nucleus or top-$k$ truncation is a \textbf{mode-sharpening} map on
the output distribution. The exact statement that matters is about a pair of tokens. At temperature
$T$ the normalizer cancels in a ratio, so for any two tokens $a, b$ at the same context:
\begin{equation}
  \log p_T(a \mid x) \;-\; \log p_T(b \mid x) \;=\; \frac{l_a(x) - l_b(x)}{T}
  \qquad \text{(exact)} .
  \label{eq:temperature}
\end{equation}

A log-odds gap is therefore scaled by exactly $1/T$ per sampling step. Below $T = 1$ that is
amplification. Retraining on those samples carries the sharpened distribution into the next
generation's weights, which is the mechanism Section~\ref{sec:collapsemap} models as a linear map in
logit space. Truncation (top-$p < 1$, top-$k$ finite) compounds it by deleting the tail outright.

Note what this does \emph{not} say: the sequence-level log-odds of two different strings does not
scale exactly by $1/T$, because the per-step normalizers no longer cancel between two different
continuations. Only the single-position decision scales exactly.

\section{Surrogates for an unobserved generation}

The threat model gives the attacker $M_{\mathrm{base}}$ and $M_0$ only. Every surrogate below
approximates $M_n$ from those two, by assuming the trajectory is locally a straight line and
extending it by a factor $n$.

\subsection{Logit space: the tilt}

\begin{equation}
  l_{\mathrm{ex}}(x) \;=\; l_{\mathrm{base}}(x) \;+\; n\left( l_0(x) - l_{\mathrm{base}}(x) \right).
  \label{eq:tilt}
\end{equation}

Computed in float32, \textbf{not clamped}. Softmax is shift invariant and is evaluated by
subtracting the row maximum, so a large logit cannot overflow; clamping to a fixed range would
instead tie every token above the ceiling and flatten everything below the floor, destroying exactly
the ranking the extrapolation produces. In distribution space this is an \textbf{exponential tilt}:
\begin{equation}
  p_{\mathrm{ex}}(\cdot \mid x) \;\propto\; p_{\mathrm{base}}(\cdot \mid x)^{\,1-n} \;\cdot\; p_0(\cdot \mid x)^{\,n},
  \label{eq:exptilt}
\end{equation}
which is the same algebra as classifier-free guidance with guidance weight $n$. Two consequences the
project uses. It is differentiable in the input embeddings through both models, so the adversarial
search can optimize against it exactly as against a real checkpoint. And it is a \emph{distribution},
not a checkpoint, so nothing is trained to produce it.

\code{extrapolate_logits} (\href{../utils/extrapolation.py}{\code{utils/extrapolation.py}}) is the
single definition, shared by the attack surrogate and the corpus generator so the two cannot drift.

\subsection{Weight space: LoRA alpha scaling}
\label{sec:weightspace}

The same first-order step taken in parameters rather than logits:
\begin{equation}
  \theta_n \;\approx\; \theta_{\mathrm{base}} \;+\; n \left( \theta_0 - \theta_{\mathrm{base}} \right).
  \label{eq:weighttilt}
\end{equation}

A LoRA layer adds $(\alpha / \mathrm{rank})\, BA$ to a frozen base weight, so the whole fine-tuning
delta is proportional to $\alpha$. Since $M_0$ is one fine-tuning step from $M_{\mathrm{base}}$,
scaling $\alpha$ by $n$ \emph{is} the weight-space counterpart of the tilt, and it costs one model
instead of two. \code{build_scaled_adapter}
(\href{../utils/extrapolation.py}{\code{utils/extrapolation.py}}).

\subsection{Data space: stacked truncation}
\label{sec:dataspace}

Collapse is driven by each generation resampling its own truncated output, so $n$ generations stack
$n$ truncations. If one step keeps the top $p_1$ of the probability mass, then
\begin{equation}
  p_n \;=\; p_1^{\,n}.
  \label{eq:topschedule}
\end{equation}

$p_1$ is \emph{fitted} to reproduce the real $M_0$, not chosen. The schedule saturates toward a
point mass as $n$ grows, which is the degenerate fixed point the real process converges to; a
schedule linear in $n$ would instead run off the end of the simplex. \code{surrogate_top_p}
(\href{../utils/extrapolation.py}{\code{utils/extrapolation.py}}).

This surrogate is a \emph{corpus-level} object. Its teacher-forced cross-entropy is identical to the
base model's, because truncation changes sampling and not likelihood, so it cannot serve as a search
proxy; the attack rejects it explicitly.

\subsection{Choosing the factor \texorpdfstring{$n$}{n}}

Three answers to one question, and they disagree, which is itself data.

\textbf{The indexing rule} $n = g + 1$ names the generation but asserts nothing about fidelity.

\textbf{The capability probe} (\code{-sf auto}) walks a ladder of factors top down and takes the
largest at which the surrogate still solves a required fraction of the clean tasks. It scans rather
than bisects, because capability is not monotonic in $n$.

\textbf{The calibration} solves, per generation, for the factor whose tilted perplexity matches the
real checkpoint's, by false position in log-perplexity:
\begin{equation}
  \text{find } n(g) \text{ such that } \quad \log \PPL_{\mathrm{ex}}\!\left( n(g) \right) \;=\; \log \PPL_{\mathrm{real}}(g).
  \label{eq:calibrate}
\end{equation}

Fitted on held-out human data, because the attacker needs the surrogate placed where generation $g$
sits on data nobody trained on. Since the fit needs checkpoints the attacker does not have, a
one-parameter law is fitted through the measured points and used to extrapolate:
\begin{equation}
  n(g) \;=\; 1 \;+\; \mathrm{scale} \cdot \ln(1 + g).
  \label{eq:law}
\end{equation}

\code{calibrated_factor} (\href{../utils/extrapolation.py}{\code{utils/extrapolation.py}}). Measured
on this repo's 0.5B run the indexing rule and the calibration diverge by many orders of magnitude in
perplexity, which is why the law exists.

\subsection{The factor the attack wants is solvable in closed form}
\label{sec:affine}

Perplexity calibration places the surrogate where generation $g$ sits as a language model. An
attack needs it to sit where generation $g$ sits \emph{at the security decision}, which is a
different quantity. That one is solvable rather than fittable.

Because the tilt is affine in the natural parameters and the partition functions cancel in a
difference of log-probabilities taken at \textbf{one} position, the surrogate's decision margin is
exactly affine in the factor:
\begin{equation}
  m(n) \;=\; (1 - n)\, m_{\mathrm{base}} \;+\; n\, m_0 .
  \label{eq:affine}
\end{equation}
Two consequences, both available from two forward passes and before any search:
\begin{equation}
  n^{*} \;=\; \frac{m_{\mathrm{base}}}{m_{\mathrm{base}} - m_0},
  \qquad
  \text{usable at all} \iff m_0 - m_{\mathrm{base}} < 0 .
  \label{eq:nstar}
\end{equation}
A non-negative slope means every increase of the factor drives the surrogate \emph{further} toward
the secure sink, so no factor whatsoever makes it model that target as insecure. This is a
guarantee from the geometry rather than a heuristic. Measured on the 0.5B, $f = 0.5$ lineage at the
pristine model's decision point, \code{fetch_user} has $n^{*} = 2.25$ and \code{load_config}
$n^{*} = 4.58$; the other three targets are unusable, \code{run_ping} because its boundary lies
behind the base model, which already prefers the shell string.

Both anchors must be scored at the \emph{same} prefix or the cancellation does not happen and the
law does not hold. \code{run_margin.py --sink} therefore evaluates every model at the pristine
model's own decision point, and \code{run_attack.py --surrogate_factor sink} reads the resulting
table (\code{sink_factor}, \href{../run_margin.py}{\code{run_margin.py}}).

\subsection{The entropy-matched tilt}
\label{sec:entropymatch}

The factor $n$ scales the \emph{difference} of the two anchors, which sets the direction of travel
but cannot set sharpness independently: as $n$ grows, the tilt~\eqref{eq:exptilt} tends to a point
mass on the argmax of the likelihood ratio $p_0 / p_{\mathrm{base}}$ and its entropy falls to zero.
The real lineage goes the other way. On the 0.5B run the mean next-token entropy
$S_M$~\eqref{eq:entropy} rises from $0.26$ nats at the base model to $0.42$ by generation~9. A tilt
large enough to model a late generation's direction is therefore far too sharp to model its
distribution, and this, not the direction, is the capability ceiling at which the surrogate stops
writing valid code. A temperature restores the missing degree of freedom.

Fix a context $x$, write $l_n = l_{\mathrm{base}} + n(l_0 - l_{\mathrm{base}})$ for the plain
tilt~\eqref{eq:tilt}, and for an inverse temperature $\beta > 0$ let
$p_\beta(l) = \softmax(\beta l)$ and $\mathbb{H}(\beta; l) = -\sum_v p_\beta(l)_v \log p_\beta(l)_v$.

\paragraph{Definition.} The entropy-matched tilt at factor $n$ is
\begin{equation}
  \tilde{l}_n(x) \;=\; \beta^{*}(x)\, l_n(x),
  \qquad
  \beta^{*}(x) \text{ the } \beta > 0 \text{ with } \mathbb{H}\!\left(\beta;\, l_n(x)\right) = \mathbb{H}\!\left(1;\, l_0(x)\right),
  \label{eq:entropymatch}
\end{equation}
equivalently $\tilde{l}_n = l_n / \tau$ with $\tau = 1/\beta^{*}$. The target is the
generation-0 anchor's entropy at the same context and under the same forbidden-token mask. It costs
no extra forward pass, since $l_0$ is already in hand for the tilt itself.

\paragraph{Existence and uniqueness.} Let $S$ be the support of the row. Then
\begin{equation}
  \frac{\partial \mathbb{H}}{\partial \beta} \;=\; -\beta \,\mathrm{Var}_{p_\beta}(l),
  \label{eq:dHdbeta}
\end{equation}
which is strictly negative unless $l$ is constant on $S$. So $\mathbb{H}$ is strictly decreasing in
$\beta$, running from $\log|S|$ as $\beta \to 0$ to $0$ as $\beta \to \infty$, and any target in
that open interval has exactly one root. The solve is a bisection on $\log\beta$ over
$[-3\ln 10,\, 3\ln 10]$, i.e.\ $\tau \in [10^{-3}, 10^{3}]$, for 12 halvings, which leaves the temperature within $0.2\%$; a bracketed bisection
on a monotone function cannot miss the root, and a target at or beyond an endpoint, such as an
exactly uniform anchor, converges to the bracket edge instead of failing. Newton with the closed
form~\eqref{eq:dHdbeta} was tried first and rejected: the variance collapses exactly where the tilt
concentrates, the step diverges, and on random logits it returned $\tau = 0.015$ where the truth
was $3.0$.

\paragraph{Rank preservation.} $\beta^{*}$ is a positive scalar per row, so it changes the order of
no pair of tokens. The argmax, and therefore greedy decoding, is identical with and without
matching. What changes is the mass behind the ranking, which is what sampling and the
cross-entropy loss see.

\paragraph{Geometry.} In the exponential family with natural parameter $\theta = l$, the plain tilt
lies on the e-geodesic through $p_{\mathrm{base}}$ and $p_0$, parameterized by $n$. Multiplying
$\theta$ by $\beta$ moves along a different e-geodesic, the ray from the uniform distribution
through $p_n$. Together the two operations span a two-parameter exponential family
\begin{equation}
  p \;\propto\; p_{\mathrm{base}}^{\,a}\; p_0^{\,b},
  \qquad (a, b) = \left(\beta(1 - n),\; \beta n\right).
  \label{eq:twoparam}
\end{equation}
The plain tilt is the slice $a + b = 1$. Entropy matching drops that constraint and replaces it with
the level-set constraint $\mathbb{H}(p) = \mathbb{H}(p_0)$: the matched surrogate is the
intersection of the ray through $p_n$ with the anchor's entropy contour. The factor chooses the
ray, the temperature chooses the point on it.

\paragraph{Interaction with the affine law.} At one decision position the partition functions
cancel in any log-odds, and a positive scalar on the logits scales that difference, so
\begin{equation}
  \tilde{m}(n) \;=\; \beta^{*}(n)\,\big[(1 - n)\, m_{\mathrm{base}} + n\, m_0\big].
  \label{eq:affinematched}
\end{equation}
Its sign is the sign of~\eqref{eq:affine}, so the boundary factor $n^{*}$ of~\eqref{eq:nstar} is
unchanged by matching; only the magnitude of the margin, and with it the scale of the hinge terms
in the objective, is rescaled. The sink-solved factor and entropy matching therefore compose
without re-solving anything.

\paragraph{Gradient.} $\tau$ is detached, so $\partial \tilde{l}_n / \partial l_n = \beta^{*}$ is
treated as a constant and the true dependence of $\beta^{*}$ on the suffix is dropped, exactly as
for a normalization constant. This is deliberate: including it would let the optimizer satisfy the
loss by shifting the sharpness rather than the direction, and the construction's point is that the
factor owns the direction and the anchor fixes the sharpness.

\paragraph{Why the anchor's entropy.} The attacker observes $p_{\mathrm{base}}$ and $p_0$ but not
the generation under attack. Since the measured entropy keeps rising past the anchor, matching to
$\mathbb{H}(p_0)$ under-estimates the late-generation entropy and the surrogate is still somewhat
too sharp, but it is a floor rather than a point mass. Extrapolating the target as well,
$\mathbb{H}^{*} = \mathbb{H}_{\mathrm{base}} + n(\mathbb{H}_0 - \mathbb{H}_{\mathrm{base}})$, is a
one-line change that has not been made, because the measured entropy trajectory is not linear in
the generation and would need the same check the calibration law~\eqref{eq:law} received.

Measured at $n = 2.25$ on the 0.5B, $f = 0.5$ lineage the plain tilt solves 1 of 5 clean tasks and
the matched tilt 4 of 5, at no measurable change in wall-clock time, since the solve touches only
logits that are already materialized. \code{entropy_matched_scale} and the \code{match_entropy}
flag of \code{extrapolate_logits} (\href{../utils/extrapolation.py}{\code{utils/extrapolation.py}});
enabled in the attack with \code{--entropy_match}.

\section{The adversarial search}

\subsection{The objective}

Let $c$ be the model under attack (the real checkpoint, or a surrogate in transfer mode) and let $H$
be the \textbf{hold set}: the models that must keep writing the secure implementation. $H$ is the
pristine baseline, plus the generation-0 anchor when it is held, so $|H|$ is 1 or 2.
\begin{align}
  L(s) \;=\;& \CE_c(w \mid s)
    && \text{(1) break} \label{eq:obj1}\\
  &+\; \nu \,\relu\!\Big( m_t - \big( \CE_c(r \mid s) - \CE_c(w \mid s) \big) \Big)
    && \text{(2) decide} \label{eq:obj2}\\
  &+\; \lambda \,\frac{1}{|H|} \sum_{h \in H} \relu\!\big( m - \CE_h(w \mid s) \big)
    && \text{(3) hold away} \label{eq:obj3}\\
  &+\; \mu \,\frac{1}{|H|} \sum_{h \in H} \CE_h(r \mid s)
    && \text{(4) hold toward} \label{eq:obj4}
\end{align}

Term by term.

\textbf{(1)} drives the insecure implementation's likelihood up on the model under attack.

\textbf{(2)} is the term that matches the loss to the verdict. (1) falling says the insecure code
got \emph{likelier}, not that it overtook the secure one, and with two models in $H$ pulling the
suffix back toward the secure code a suffix can improve (1) indefinitely without ever flipping the
answer. (2) is a hinge on the target's own decision margin and goes slack once the insecure code
leads by $m_t$ nats per token. $m_t = 0$ disables it and recovers the previous objective.

\textbf{(3)} is a hinge, not gradient ascent: unbounded ascent on $\CE_h(w)$ diverges. It stops
contributing once a hold model is already $m$ nats away from the insecure code.

\textbf{(4)} anchors the hold models on the secure implementation, so that ``did not emit the
insecure string'' cannot be satisfied by making them emit garbage instead.

\textbf{The $1/|H|$ is load-bearing.} Summing instead of averaging doubles the pressure to hold the
moment the anchor is added, while the single break term stays as it is, so $\lambda$, $\mu$ and $m$
silently stop meaning what they were calibrated to mean. Measured on one cell at a fixed budget, the
summed objective reached a best total of $5.90$ and found no hit where the averaged one hits at
$3.39$. \code{run_selective_attack.py} generalizes the same rule to an arbitrary hold group.

\code{ContrastiveGCG.objective} and \code{.combined_gradient}
(\href{../run_attack.py}{\code{run_attack.py}}).

\subsection{The sink-conditional objective}
\label{sec:sinkobj}

Every cross-entropy in \eqref{eq:obj1}--\eqref{eq:obj4} is a mean over the whole reference
implementation, a hundred-odd tokens of which two or three decide the verdict. Measured on the
0.5B, $f = 0.5$ lineage that mean fell from about $3$ to about $0.8$ nats per token while the
surrogate's own output was unclassifiable in 75 of 75 checks: the suffix had pushed the model off
the task rather than across the decision.

Let $k$ be the length of the shared token prefix of $r$ and $w$, and let $r_k$ and $w_k$ be the
first differing spans of the two references (a pure insertion is extended by one shared token so
both are non-empty). For any model $M$ define
\begin{align}
  \mathrm{pref}_M(s)   \;&=\; \log p_M(w_k \mid s, r_{<k}) \;-\; \log p_M(r_k \mid s, r_{<k}),
  \label{eq:pref}\\
  \mathrm{prefix}_M(s) \;&=\; \frac{1}{k} \sum_{t < k} -\log p_M(r_t \mid s, r_{<t}).
  \label{eq:prefixce}
\end{align}
Both spans are teacher-forced on the same prefix, so \eqref{eq:pref} is a difference of two
log-probabilities at one context: a decision margin, the quantity the static classifier reads and
the quantity Section~\ref{sec:affine} measures. The sink-conditional objective is
\begin{align}
  L_{\mathrm{sink}}(s) \;=\;& \zeta\, \relu\!\big( m_s - \mathrm{pref}_c(s) \big)
    && \text{target prefers the insecure span} \label{eq:sink1}\\
  &+\; \kappa\, \mathrm{prefix}_c(s)
    && \text{and still writes the prefix} \label{eq:sink2}\\
  &+\; \lambda\, \frac{1}{|H|} \sum_{h \in H} \relu\!\big( m + \mathrm{pref}_h(s) \big)
    && \text{each hold model prefers the secure span} \label{eq:sink3}\\
  &+\; \mu\, \frac{1}{|H|} \sum_{h \in H} \CE_h(r \mid s)
    && \text{and keeps writing the secure code.} \label{eq:sink4}
\end{align}

Three properties. Under the logit surrogate $\mathrm{pref}_c$ is exactly affine in the factor, by
the argument of Section~\ref{sec:affine}, so each target's boundary $n^{*}$ is the natural factor
to search at, and \code{--surrogate_factor sink} now moves the surrogate to each target's own root
(plus \code{--sink_offset}) when its search starts. The entropy-matching temperature of
Section~\ref{sec:entropymatch} is a positive scalar on the logits and scales $\mathrm{pref}$
without changing its sign. And $m_s$ is in nats of a two-way log-odds, so it maps directly onto
the sampled verdict: at temperature $T$ the insecure span wins the two-way race with probability
$\sigmoid(m_s / T)$, about $95\%$ at $m_s = 2$, $T = 0.7$.

The full-sequence means are still computed and recorded, so \code{--max_wrong_loss} keeps its
meaning. \code{ContrastiveGCG._sink_objective} and \code{._sink_gradient}
(\href{../run_attack.py}{\code{run_attack.py}}); result files carry \code{_sinkobj}.

\subsection{The discrete step}

The suffix is represented as a one-hot matrix $X \in \{0,1\}^{|s| \times V}$, and the input
embeddings are $e = XE$ for embedding matrix $E$. The loss is differentiable in $X$, giving
\begin{equation}
  G \;=\; \nabla_X L(s) \;\in\; \mathbb{R}^{|s| \times V}.
  \label{eq:grad}
\end{equation}

$G_{i,v}$ is a first-order estimate of the change in $L$ from putting token $v$ at position $i$. One
GCG step is then:
\begin{enumerate}
  \item for each position $i$, take the $k$ tokens with the most negative $G_{i,\cdot}$ as
    candidates (implemented as the top-$k$ of $-G$, with disallowed tokens set to $+\infty$ first so
    they can never be selected, which is how the non-ASCII and special-token constraints enter);
  \item draw $B$ candidate suffixes, each replacing $n_{\mathrm{replace}}$ positions uniformly from
    those sets;
  \item discard candidates that do not survive a detokenize/retokenize round trip;
  \item evaluate $L$ \textbf{exactly} on all survivors and keep the $\argmin$.
\end{enumerate}

The gradient only proposes; the exact evaluation decides. The search is greedy coordinate descent in
this sense and never accepts a worse candidate. Steps 1 and 3 are \code{sample_ids_from_grad} and
\code{filter_ids} (\href{../utils/gcg.py}{\code{utils/gcg.py}}).

Note that this makes the search itself \textbf{decoding-blind}: teacher-forced cross-entropy
involves no sampling, which is why the data-space surrogate of Section~\ref{sec:dataspace} is
useless here and why the verdict of Section~\ref{sec:verdict} has to be imposed separately.

\section{The verdict}
\label{sec:verdict}

\subsection{The hit rule}

Let $\mathrm{status}_M \in \{\text{secure}, \text{insecure}, \text{unclassified}\}$ be the static
classification of what model $M$ emits for the adversarial prompt. A hit requires all three of:
\begin{align}
  \mathrm{status}_c &= \text{insecure} && \text{the model under attack writes the vulnerability} \\
  \mathrm{status}_b &= \text{secure} && \text{the pristine baseline does not} \\
  \mathrm{status}_a &= \text{secure} && \text{the generation-0 anchor does not \quad (when held)}
\end{align}

The third condition is what makes the claim about \emph{collapse} rather than about fine-tuning:
without it, a suffix that breaks every descendant of the base model counts, and that is a statement
about the first fine-tune, not about the trajectory.

\subsection{Sampled majority}

Each model is decoded $k$ times at the deployment settings and charged with the \textbf{majority}
status, ties resolved against the hit. If a single sample is insecure with probability $q$, the
model is charged insecure iff
\begin{equation}
  \Bin(k, q) \;>\; k/2 .
  \label{eq:majority}
\end{equation}

So the sampled verdict is, to first order, a threshold at $q = 1/2$, where the greedy verdict is a
threshold on the argmax. \textbf{The two verdicts are two thresholds on one underlying quantity},
which is the subject of Section~\ref{sec:margin} and the reason the project measures that quantity
directly.

\subsection{Loss and survival}

Under ancestral sampling the probability of emitting the target exactly is
\begin{equation}
  P(y) \;=\; \exp\!\left( -|y| \cdot \CE(y) \right).
  \label{eq:survival}
\end{equation}

So a target of 40 tokens at $\CE = 0.05$ nats/token is emitted about 14\% of the time per draw, and
at $\CE = 0.02$ about 45\%. This is what \code{--max_wrong_loss} thresholds: it refuses to count a
hit until the insecure code is likely enough to survive being sampled, instead of accepting the
first argmax flip. Truncation raises these numbers, since mass deleted from the tail is
redistributed.

\section{The security margin}
\label{sec:margin}

A decoding-independent scalar that both verdicts are thresholds on. Measured on the \textbf{clean}
prompt to characterize a model, and on the adversarial prompt to characterize a suffix.
Implemented in \code{measure()} (\href{../run_margin.py}{\code{run_margin.py}}).

\subsection{Three quantities}

Let $k$ be the first position at which the tokenizations of $r$ and $w$ diverge, with the shared
prefix teacher-forced.

\textbf{Decision margin}, the log-odds of the secure token over the insecure one at the branch
point:
\begin{equation}
  \mathrm{dm}_M \;=\; \log p_M\!\left( r_k \mid \text{prefix} \right) \;-\; \log p_M\!\left( w_k \mid \text{prefix} \right).
  \label{eq:dm}
\end{equation}

\textbf{Sequence margin}, the same at the level of whole implementations:
\begin{equation}
  \mathrm{sm}_M \;=\; \log p_M(r) - \log p_M(w) \;=\; |w| \cdot \CE_M(w) \;-\; |r| \cdot \CE_M(r).
  \label{eq:sm}
\end{equation}

\textbf{Entropy}, the mean next-token entropy over the secure target's positions, which is a direct
read-out of the sharpening of Section~\ref{sec:sharpening} and serves as the collapse-severity axis:
\begin{equation}
  S_M \;=\; \frac{1}{|r|} \sum_t \mathbb{H}\!\left( p_M(\cdot \mid r_{<t}) \right).
  \label{eq:entropy}
\end{equation}

\subsection{How the margins relate to the verdicts}

Greedy decoding emits the secure token at the branch point iff it is the argmax, which requires
$\mathrm{dm}_M > 0$ but is not implied by it, since a third token may lead. The measurement records
\code{secure_is_argmax} separately for exactly this reason, and it is often false: the model phrases
the function differently from both references, in which case $\mathrm{dm}$ compares two tokens the
model was never going to write and $\mathrm{sm}$ is the quantity to read.

By the exact result of Section~\ref{sec:sharpening}, the decision margin at temperature $T$ is
$\mathrm{dm}/T$. So lowering the temperature scales the decision margin up and makes a flip harder,
and raising it scales toward zero and makes the outcome a coin flip. This is the formal content of
``exploitability is decoding-conditional''.

\subsection{Headroom, and its limits}

For the three-sided hit rule to be satisfiable, the model under attack must be \emph{closer} to the
insecure variant than the anchor it has to spare. Hence
\begin{equation}
  \mathrm{headroom}(g, \mathrm{task}) \;=\; \mathrm{sm}_{M_0}(\mathrm{task}) \;-\; \mathrm{sm}_{M_g}(\mathrm{task}).
  \label{eq:headroom}
\end{equation}

Positive headroom means collapse has already moved the target in the attack's favour relative to the
model that must be held. Negative headroom means the search is asked to drag the target past a model
it must simultaneously spare.

\textbf{Headroom is a prior, not a predictor.} Measured on the $f = 0.5$ lineage, the only cell that
produced hits was ranked last of its peers by both the sequence and the decision form. A
clean-prompt margin says where a model starts, not how far a suffix can move it; the quantity that
would predict success is a sensitivity, not a level. Use the ranking to report and to diagnose, not
to filter.

\subsection{The admissibility band}

Two conditions must hold at once for a cell $(g, \mathrm{task})$ to support the claim. The
capability gate requires $M_g$ to write \emph{secure} code unaided, and headroom requires $M_g$ to
be \emph{closer to insecure} than $M_0$. These pull in opposite directions, and their intersection
is the real object of study. Measured on the $f = 0.5$ lineage over 45 cells:

\begin{center}
\begin{tabular}{lr}
\toprule
cells & count \\
\midrule
total & 45 \\
pass the capability gate & 14 \\
positive headroom & 13 \\
both & 3 \\
\bottomrule
\end{tabular}
\end{center}

The clearest illustration is the OS-command-injection target, which has the largest headroom at
every generation and is inadmissible at every generation because the anchor already writes the
vulnerability with no suffix at all.

\section{The collapse map (poisoning)}
\label{sec:collapsemap}

The poisoning line models collapse as a map on a scalar \textbf{observable} rather than on weights.
Let $r_n$ be the trigger-context payload rate at generation $n$: the fraction of answers to
trigger-bearing instructions that contain the payload. One generation resamples and retrains, which
by Section~\ref{sec:sharpening} is near-linear in logit coordinates:
\begin{equation}
  \logit(r_{n+1}) \;=\; a \cdot \logit(r_n) \;+\; b .
  \label{eq:map}
\end{equation}

with $a$ the sharpening slope. Section~\ref{sec:sharpening} predicts $a \approx 1/T$, about $1.43$ at
$T = 0.7$, but $a$ is \emph{measured}, never assumed. The map is affine, so it has a single fixed
point
\begin{equation}
  L^{*} \;=\; \frac{b}{1-a}, \qquad r^{*} \;=\; \sigmoid(L^{*}).
  \label{eq:fixedpoint}
\end{equation}

In the shifted coordinate $u_n = \logit(r_n) - L^{*}$ the map is pure geometric growth, $u_n = a^n
u_0$, giving three closed forms:
\begin{align}
  \text{forecast} && r_n &\;=\; \sigmoid\!\left( L^{*} + a^{n} \left( \logit(r_0) - L^{*} \right) \right) \label{eq:forecast}\\
  \text{activation} && n^{*} &\;=\; \frac{\log\!\left( -L^{*} / \left( \logit(r_0) - L^{*} \right) \right)}{\log a} \label{eq:activation}\\
  \text{required dose} && r_0 &\;=\; \sigmoid\!\left( L^{*} \left( 1 - a^{-n^{*}} \right) \right) \label{eq:dose}
\end{align}
where activation is the generation at which greedy decoding flips, i.e.\ $\logit(r) = 0$.

\textbf{The fixed point is unstable iff $a > 1$}, and a dormant payload amplifies iff additionally
$r^{*} < 0.5$, so that the initial rate sits on the growing side. When those fail, no generation-0
dose arms the fuse and the tools return infeasible rather than a number. That is a finding about the
attack surface, not a tuning failure: the LoRA fine-tune damps the sharpening.

Implemented in \href{../utils/collapse_map.py}{\code{utils/collapse_map.py}}, derived in
\href{collapse_map_poison.md}{\code{collapse_map_poison.md}}.

The earlier gradient-matching variant optimizes poison content against the \emph{weight} surrogate
of Section~\ref{sec:weightspace}. If a reference backdoor sits at $\theta^{*}$, then under
$\theta_k \approx \theta_{\mathrm{base}} + k(\theta_0 - \theta_{\mathrm{base}})$ the generation-0
delta needed to arrive there at generation $k$ is
\begin{equation}
  \Delta W_0 \;=\; \frac{\Delta W^{*}}{k},
  \label{eq:gradmatch}
\end{equation}
and the poison text is searched to make the training gradient point at it, by minimizing
$-\cos\!\left( \nabla_B L,\, B^{*} \right)$ with the LoRA basis frozen. See
\href{gradient_matching_poison.md}{\code{gradient_matching_poison.md}}. It assumes weight-space
collapse, which is not the regime this pipeline is in, which is why the data-space map above
replaced it.

\section{Symbol table}

\begin{center}
\begin{tabular}{lll}
\toprule
symbol & meaning & default \\
\midrule
$n$ & surrogate extrapolation factor & $g + 1$ \\
$n^{*}$ & factor at which the surrogate's decision margin is 0, eq.~\eqref{eq:nstar} & solved \\
$\beta^{*}$, $\tau$ & entropy-matching inverse temperature and temperature, eq.~\eqref{eq:entropymatch} & solved \\
$f$ & real-data fraction in the training mixture & 0 \\
$T$, $p$, $k$ & sampling temperature, nucleus, top-$k$ & 0.7, 0.8, 20 \\
$\lambda$ & weight of the hold-away hinge & 1.0 \\
$\mu$ & weight of the hold-toward term & 0.5 \\
$m$ & hold-away hinge margin & 3.0 \\
$\nu$, $m_t$ & weight and margin of the target-side hinge & 1.0, 0 (off) \\
$\zeta$, $m_s$, $\kappa$ & sink-conditional break weight, margin, stay-on-task weight & 1.0, 2.0, 1.0 \\
$B$, $k_{\mathrm{top}}$ & GCG search width, gradient top-$k$ & 256, 256 \\
$H$ & hold set, $\{\text{baseline}\}$ or $\{\text{baseline}, \text{anchor}\}$ & with anchor \\
$a$, $b$, $L^{*}$ & collapse-map slope, intercept, fixed point & measured \\
$r_n$ & trigger-context payload rate at generation $n$ & measured \\
\bottomrule
\end{tabular}
\end{center}

\section{Where each formula lives}

\begin{center}
\begin{tabular}{ll}
\toprule
formula & file \\
\midrule
per-token cross-entropy, objective, gradient & \href{../run_attack.py}{\code{run_attack.py}} \\
sink-conditional objective, sink spans & \href{../run_attack.py}{\code{run_attack.py}} \\
GCG candidate step & \href{../utils/gcg.py}{\code{utils/gcg.py}} \\
logit tilt, alpha scaling, top-$p$ schedule, calibration law & \href{../utils/extrapolation.py}{\code{utils/extrapolation.py}} \\
entropy-matched tilt & \href{../utils/extrapolation.py}{\code{utils/extrapolation.py}} \\
affine law, boundary factor & \href{../run_margin.py}{\code{run_margin.py}} \\
data mixture & \href{../run_baseline.py}{\code{run_baseline.py}} \\
margins, entropy, headroom & \href{../run_margin.py}{\code{run_margin.py}} \\
sampled majority verdict & \href{../run_attack.py}{\code{run_attack.py}} \\
collapse map and its closed forms & \href{../utils/collapse_map.py}{\code{utils/collapse_map.py}} \\
\bottomrule
\end{tabular}
\end{center}

\end{document}
